\documentclass[UTF8,fontset=fandol,11pt]{ctexart}
\usepackage[a4paper,margin=24mm]{geometry}
\usepackage{amsmath,amssymb,amsthm,graphicx,tabularx,booktabs,hyperref}
\newtheorem{definition}{定义}[section]
\newtheorem{theorem}[definition]{定理}
\newtheorem{example}[definition]{例题}
\newtheorem{exercise}[definition]{练习}
\hypersetup{unicode=true,colorlinks=true,linkcolor=blue}
\title{答案定位}
\date{}
\begin{document}
\maketitle
\par\noindent 离线合成演示：不代表真实模型对本专题的教学验收。\par
\section{第一节}\label{sec:S01}
% UNIT EX
\begin{example}
计算 $1+1$。
\end{example}
% UNIT EX.A
\par\noindent\textbf{解：}例题原解 $1+1=2$。\par
% UNIT E1
\begin{exercise}
\label{mathagent:exercise:d038ac307b6079d71d636e7dfa039fc3fa1e094f72c4fea758a79f0f0081f728}
计算 $2+2$。
\end{exercise}
\section{第二节}\label{sec:S02}
% UNIT E1
\begin{exercise}
\label{mathagent:exercise:84928b6d72a05fc19d6c10c013040b9e7d1f9cf2265eec91ca520a387ed1049f}
计算 $3+3$。
\end{exercise}
\section*{练习答案}
% UNIT A1
\par\noindent\textbf{\hyperref[mathagent:exercise:d038ac307b6079d71d636e7dfa039fc3fa1e094f72c4fea758a79f0f0081f728]{练习 \ref*{mathagent:exercise:d038ac307b6079d71d636e7dfa039fc3fa1e094f72c4fea758a79f0f0081f728} 的解答：}}第一节答案 $2+2=4$。\par
% UNIT A1
\par\noindent\textbf{\hyperref[mathagent:exercise:84928b6d72a05fc19d6c10c013040b9e7d1f9cf2265eec91ca520a387ed1049f]{练习 \ref*{mathagent:exercise:84928b6d72a05fc19d6c10c013040b9e7d1f9cf2265eec91ca520a387ed1049f} 的解答：}}第二节答案 $3+3=6$。\par
\end{document}
